% -----------------------------------------------------------------------------
% Section 7: Conclusion
% Paper: "Silent Decisions Are Type Errors: Enforcing AI Accountability
%         via Linear Resource Types"
% Venue: IEEE Computer — Special Issue on AI Governance
% Deadline: 13 April 2026
% Authors: Daniel Lau Pereira Soares
% -----------------------------------------------------------------------------

\section{Conclusion}
\label{sec:conclusion}

This paper began from a simple observation: regulatory frameworks
such as LGPD Art.~18\S2 and EU AI Act Art.~86 mandate that
explanations and contestability rights \emph{accompany} a decision,
but they cannot prevent a decision from existing \emph{without} them.
We called this failure mode a \emph{silent decision}, and argued
that its root cause is not a failure of intent but a failure of
type---a structural property of architectures that treat
accountability as a runtime audit obligation rather than a
compile-time invariant.

Drawing on Girard's Linear Logic~\cite{girard1987linear} and Rust's
ownership model~\cite{jung2017rustbelt}, we introduced the
Bound-To-Verdict (BTV) framework and its core type law:
$V \multimap (E \otimes C)$.
Under this invariant, a \texttt{Verdict} cannot be constructed
without atomically consuming an \texttt{EvidenceToken} and a
\texttt{ComplianceToken}.
The three structural protections---private fields, \texttt{pub(crate)}
consume, and \texttt{\#[must\_use]}---close every syntactic path
around the constructor, making a silent decision a type error
detected at compile time.

We proved the \emph{Constitutional Enclosure Theorem}: in any Safe
Rust program that includes the BTV crate, the set $\mathcal{S}$ of
materialized silent decisions is empty by construction.
The proof is exhaustive (three cases), machine-checkable
(\texttt{cargo test}: six clauses, four seconds), and accompanied
by a Fail-Secure Corollary: a system that cannot produce evidence
cannot issue a verdict---it fails securely rather than silently.

Empirical benchmarks on the reference implementation confirm that
this static guarantee is a zero-cost abstraction in runtime.
The full governance cycle for a 4\,KB payload completes in
$1.67\,\mu$s; every nanosecond of that cost is attributable to
cryptographic primitives (BLAKE3 + HMAC-SHA256) that any
accountable system must execute regardless of architecture.
The type-level enforcement itself contributes zero overhead.

The implications extend beyond compliance engineering.
If accountability can be expressed as a type invariant, then
the correctness of governance becomes a property that compilers
can check, proof assistants can verify, and regulators can audit
through \texttt{cargo test} rather than through manual code
review.
Silent decisions are not an inevitable cost of deploying AI at
scale---they are a bug that a sufficiently expressive type system
eliminates entirely.
